%
% CVPR preamble.
%
% cvpr.sty already provides: times, xspace, xcolor[dvipsnames], graphicx,
% amsmath, amssymb, booktabs, natbib[numbers,sort&compress], silence, etoolbox,
% caption, subcaption, url and enumitem -- and it loads cleveref itself via
% \AtEndPreamble (i.e. strictly after hyperref), so none of those may be loaded
% here. (This is the same package set wacv.sty supplied; the ICLR port had to
% declare them all by hand because iclr2027_conference.sty supplied only
% fancyhdr + natbib.)
%
% Disable PDF object-stream compression: the large architecture figure this
% was added for (the former imgs/GenAU_Arch_final.pdf) otherwise overflowed pdfTeX's
% object-stream buffer ("TeX capacity exceeded [PDF object stream buffer]").
\ifdefined\pdfobjcompresslevel \pdfobjcompresslevel=0 \fi

% --- required by this paper and NOT provided by cvpr.sty ---
\usepackage{multirow}   % \multirow in the result/ablation tables
\usepackage{colortbl}   % \rowcolor{gray!20} (the "table" feature of xcolor)
\usepackage{tabularx}   % column-width ablation table (Table 2)
\usepackage{tikz}       % motivation radar chart (Fig. 1)
\usepackage{microtype}  % char protrusion + font expansion: tighter lines, fewer runts

% --- compact (but standard) float spacing to help fit the page limit ---
\setlength{\textfloatsep}{10pt plus 2pt minus 2pt}      % top/bottom single-col floats
\setlength{\intextsep}{10pt plus 2pt minus 2pt}         % [h] floats
\setlength{\dbltextfloatsep}{12pt plus 2pt minus 2pt}   % full-width figure*/table*

% --- ablation-table deltas to the default row: fixed-width box keeps the numbers aligned
\newcommand{\abd}[2]{\makebox[1.9em][l]{\textsubscript{\textcolor{#1}{#2}}}}
\newcommand{\abup}[1]{\abd{ForestGreen}{$+$#1}}
\newcommand{\abdn}[1]{\abd{BrickRed}{$-$#1}}
\newcommand{\abeq}{\abd{gray}{$\pm$0.0}}
\newcommand{\abno}{\abd{black}{}}

% --- inline annotations
\newcommand{\red}[1]{{\color{red}#1}}
\newcommand{\todo}[1]{{\color{red}#1}}
\newcommand{\TODO}[1]{\textbf{\color{red}[TODO: #1]}}
% --- disable by uncommenting
% \renewcommand{\TODO}[1]{}
% \renewcommand{\todo}[1]{#1}
